\chapter{Experimental Design}
\label{chap:design}

The experiment uses the simulation-design framework of \citet{morris2019}
to evaluate whether Expanded Welch's larger rejection cutoff improves
performance compared with Normal Wald. It first measures false
positives when the two populations have equal MI, then examines the
estimation errors behind those results,
and finally estimates power when MI differs between the two populations.
Bias and variance diagnostics
link rejection rates to errors in the estimated MI difference and standard
error. Each comparison starts from fixed joint probability tables with
known MI values, so the
true difference is available when assessing the sampled estimates.

\section{Questions and controlled comparisons}

The first study measures false positives when the two populations have
equal MI. It varies sample size, table size and category frequencies to
determine when Normal Wald and Expanded Welch reject equal MI too often
or too rarely relative to the intended 5\% rate. It also records how often
each test returns a usable result.

The second study examines how the size of the MI values affects Expanded
Welch's estimated degrees of freedom. It changes the MI value in both
populations while keeping their MI values equal and holding table size
and category frequencies fixed. The study includes small positive MI values
to examine how the degrees of freedom behave close to independence.

The third study measures statistical power when MI differs between the
two populations. It keeps one population fixed and increases the MI of the
other, while holding sample size, table size and category frequencies fixed
within each comparison. Comparing Normal Wald and Expanded Welch shows
how the larger rejection cutoff affects the probability of detecting a
true MI difference.

Two supporting checks examine how population construction and sample
allocation affect false-positive rates and power. The construction check
compares two ways of introducing dependence between \(X\) and \(Y\), with
population MI values, category frequencies and sample sizes matched across
the two constructions. The sample-allocation check keeps both populations
and the total sample size fixed but varies how many observations are drawn
from each population. Table~\ref{tab:design-families} summarises these
studies and the construction and sample-allocation checks.

\begin{table}[H]
\centering
\caption{Overview of the three main studies and two supporting checks.}
\label{tab:design-families}
\begin{tabularx}{\textwidth}{@{}l>{\raggedright\arraybackslash}X>{\raggedright\arraybackslash}X@{}}
\toprule
Comparison & Factor varied & Main purpose\\
\midrule
Equal MI & Sample size, table size and category frequencies & False positives and usable results.\\
Population MI & MI value, equal in both populations & Degrees of freedom and near-independence behaviour.\\
Power & True MI difference & Power alongside null calibration.\\
\midrule
Construction check & Copying or four-cell transfer & Sensitivity to a different association pattern.\\
Sample allocation & Relative sample sizes at fixed total & False positives and power with unequal samples.\\
\bottomrule
\end{tabularx}
\end{table}

\section{Constructing joint distributions with known MI}

The construction changes association while preserving the category
frequencies. For one population, draw \(X\) with category probabilities
\(a=(a_1,\ldots,a_k)\). With probability \(\lambda\), set \(Y=X\);
otherwise draw \(Y\) independently with the same probabilities \(a\).
Copying makes matching categories more likely than under independence.
Both branches give \(Y\) the same marginal probabilities, so changing
\(\lambda\) does not change the row or column totals.

An observation in a diagonal cell can arrive through either branch,
whereas an off-diagonal observation can arrive only through the
independent branch.
Combining these probabilities gives
\begin{equation}
 p_{ij}(\lambda)=(1-\lambda)a_i a_j
                  +\lambda a_i\mathbf{1}_{\{i=j\}},
 \qquad 0\leq\lambda<1,
 \label{eq:copying-construction}
\end{equation}
where the indicator is one when the two categories match and zero otherwise.
This construction is a special case of the mixture model described by
\citet{bocci2010}.
At \(\lambda=0\), the variables are independent and MI is zero.
Increasing \(\lambda\) increases MI, so the target MI determines the required
copying probability. Appendix~\ref{app:construction} gives the numerical
search and shows why the margins remain fixed.

\begin{samepage}
\medskip
\thesisminiheading{Constructing a \(2\times2\) joint distribution with positive MI.}

For \(P\), take \(a=(0.7,0.3)\) and a copying probability of
approximately 0.20347. Cell \((1,1)\) requires \(X=1\), which has
probability 0.7. Given \(X=1\), copying always gives \(Y=1\), whereas
an independent draw gives \(Y=1\) with probability 0.7. Thus
\begin{align*}
 p_{11}&\simeq0.7\{0.20347+(1-0.20347)(0.7)\}=0.53273,\\
 p_{12}&\simeq0.7(1-0.20347)(0.3)=0.16727.
\end{align*}
\end{samepage}
The second equation has no copying term because copying cannot produce
\(X=1,Y=2\). The remaining cells follow the same rule.
For \(Q\), use \(a=(0.8,0.2)\) and a copying probability of approximately
0.26091. These choices give \(I(P)=0.02\) and \(I(Q)=0.03\) nats:
\begin{equation}
 P\simeq\begin{pmatrix}0.53273&0.16727\\0.16727&0.13273\end{pmatrix},
 \qquad
 Q\simeq\begin{pmatrix}0.68175&0.11825\\0.11825&0.08175\end{pmatrix}.
 \label{eq:design-example-pair}
\end{equation}
Both joint distributions have positive MI, and the true difference used
to assess power has magnitude 0.01 nats.

\medskip
\thesisminiheading{Table sizes and margin regimes.}

The copying construction is used for square tables of sizes
\(2\times2,4\times4,8\times8\) and \(16\times16\), containing
4, 16, 64 and 256 cells. Each doubling of the number of categories
multiplies the number of cells by four, so at a fixed sample size the
observations are distributed over more possible cells.

Three margin regimes vary the category frequencies: uniform frequencies
in both populations, the same uneven frequencies in both, or different
uneven frequencies between them. For uneven margins, the first category
has probability \(\rho\), and the remaining mass is shared equally:
\begin{equation}
 m_k(\rho)=\left(\rho,\frac{1-\rho}{k-1},\ldots,
                   \frac{1-\rho}{k-1}\right).
 \label{eq:skew-margin}
\end{equation}
Each population uses the same vector for its row and column margins.
Table~\ref{tab:margin-profiles} specifies the vectors for \(P\) and \(Q\).
The copying probabilities for \(P\) and \(Q\) are chosen separately to
achieve the targets \(I(P)\) and \(I(Q)\).

\begin{table}[htbp]
\centering
\caption{Row and column margins for each population.}
\label{tab:margin-profiles}
\begin{tabular}{@{}lll@{}}
\toprule
Regime & \(P\): rows and columns & \(Q\): rows and columns\\
\midrule
Uniform & \((1/k,\ldots,1/k)\) & \((1/k,\ldots,1/k)\)\\
Same skew & \(m_k(0.8)\) & \(m_k(0.8)\)\\
Different skew & \(m_k(0.7)\) & \(m_k(0.8)\)\\
\bottomrule
\end{tabular}
\end{table}

Under equal MI, the uniform and same-skew regimes give identical joint
probability tables, \(P=Q\). In contrast, the different-skew regime gives
different tables with equal MI because the category frequencies differ
between populations. Comparing the regimes therefore tests false-positive
behaviour under both identical and distinct population tables.

Same skew provides a matched baseline for different skew in all three
studies and both supporting checks. With equal MI and equal sample sizes,
the same-skew samples have identical distributions, so their expected MI
estimation biases cancel in the difference. Different-skew populations can
retain unequal biases despite having equal MI. Under an MI difference or
unequal sample sizes, matching margins alone does not ensure that biases
cancel. The two regimes therefore distinguish failures that also occur
with matching category frequencies from those observed when the frequencies
differ.

Increasing table size under the uneven-margin rule spreads the remaining
mass over more categories, changing both the number of cells and their
probabilities. This comparison shows what happens under a common margin
rule as the table grows. Uniform versus same skew instead compares even
and uneven category frequencies at a fixed size; different skew allows the
two populations' frequencies to differ as well.

\section{The three studies and supporting checks}

\subsection{Sample size, table size and margins}

The first study fixes \(I(P)=I(Q)=0.02\) nats and crosses all four sizes
with all three margin regimes. For each of the twelve fixed population
pairs, sample size varies over
\begin{equation}
 n_P=n_Q\in\{10,20,50,100,250,500,1000,2500,10000,50000\}.
 \label{eq:null-sample-grid}
\end{equation}
These 120 settings show whether the false-positive rate approaches 5\%
as more observations become available, and whether table size or margins
change that progression. Valid-result rates distinguish a genuinely low
false-positive rate from a failure to calculate the test.

The common target of 0.02 nats gives positive MI in all four table sizes
and three margin regimes. Fixing population MI does not fix the
difficulty of estimating it: tables with equal MI can have different
cell probabilities and different MI sampling variances. Keeping each
population pair fixed throughout the sweep isolates the effect of sample
size on test performance.
The sweep includes small samples, intermediate samples and large-sample
checks so that both calculation failures and convergence are visible.

\subsection{Population MI and degrees of freedom}

The second study uses \(4\times4\) tables under both same skew and
different skew. Within each regime, it keeps the margins fixed and varies
the positive MI value while keeping it equal in both populations:
\begin{equation}
 I(P)=I(Q)\in\{0.0001,0.0003,0.001,0.003,0.01,0.02,0.05\}.
 \label{eq:association-grid}
\end{equation}
The sweep is repeated at \(n=100\) and 1,000. Small positive MI targets
approach independence, where the first-order MI variance becomes small,
while larger targets show how the tests behave away from that region. Tracking
the estimated degrees of freedom connects changes in rejection to the
cutoff from the Student \(t\) distribution. The two sample sizes show
whether these patterns persist with more observations. All targets are
positive, so the sweep approaches rather than reaches exact independence.
Comparing the two regimes shows whether the relationship between population
MI and estimated degrees of freedom persists when category frequencies
match. The two regimes and two sample sizes give 28 settings, four of which
are already in the first study.

\subsection{Power to detect MI differences}

The third study fixes \(I(P)=0.02\) and sets \(I(Q)=0.02+|\Delta|\),
using
\begin{equation}
 |\Delta|\in\{0,0.001,0.002,0.005,0.01,0.02,0.05,0.1,0.2\}
 \quad\text{nats}.
 \label{eq:detection-grid}
\end{equation}
At \(n=1{,}000\) and 10,000, compare all four sizes under both same skew
and different skew, with an additional uniform-margin comparison at
\(4\times4\).
The signed parameter is \(\Delta=I(P)-I(Q)\), which is negative
when \(Q\) has the larger MI; figures use its known nonnegative magnitude.

Within each curve, the entire population table \(P\) stays fixed, while
\(Q\)'s copying probability increases to give the specified MI differences.
Small differences show where the methods may
disagree; larger differences show whether power eventually becomes high.
Each curve includes the corresponding equal-MI setting from the null
study, so power can be read alongside the false-positive rate at the
same sample size.

\subsection{A supporting change in construction}

This check examines whether the comparison between Normal Wald and Expanded
Welch changes when dependence is distributed differently across the
population tables. It compares the copying construction with a four-cell
construction, keeping table size, category frequencies, population MI and
sample sizes matched between constructions.

The four-cell construction starts from a table in which \(X\) and \(Y\)
are independent. It adds the same amount of probability to cells \((1,1)\)
and \((2,2)\) and subtracts it from \((1,2)\) and \((2,1)\), preserving
all row and column totals. Unlike copying, which changes probabilities
throughout the table, this transfer changes only four cells. It follows
the sparse construction of \citet[Section~3(b)]{berrett2021};
Appendix~\ref{app:construction} gives its formula and target-MI search.

The check uses \(4\times4\) and \(16\times16\) tables under both same
skew and different skew, with \(n=1{,}000\) and 10,000 observations per
population. The MI pairs \((I(P),I(Q))=(0.02,0.02)\) and
\((0.02,0.03)\) nats give the false-positive and power comparisons,
respectively. These choices give sixteen four-cell settings, each matched
to a copying setting from the main studies. The comparison measures
rejection rates, valid-result frequencies, estimation errors and the
estimated direction when population MI differs.

\subsection{A supporting unequal-sample comparison}

This check examines how dividing a fixed total of 2,000 observations
unequally between the two populations affects false-positive rates and
statistical power.

For each fixed population pair, it compares
\((n_P,n_Q)=(1000,1000),(400,1600)\) and \((1600,400)\).
The first allocation gives the populations equal sample sizes; the other
two give one population four times as many observations as the other.
Only the sample sizes change: the population tables and their MI values
remain fixed.

The check uses \(4\times4\) copying populations under both same skew and
different skew. In each regime, \(I(P)=I(Q)=0.02\) nats gives the
false-positive comparison, while \(I(P)=0.02\) and \(I(Q)=0.03\) nats
gives the power comparison. Both tests use their existing estimators and
validity rules for all three allocations. The four balanced settings are
already included in the main studies, so the unequal allocations add eight
settings.

\medskip
\thesisminiheading{A supporting bias-correction check.}

A jackknife check changes the MI bias correction while keeping the
original standard errors and standard normal distribution cutoff,
using the same samples from four existing \(16\times16\) different-skew
copying settings. Appendix~\ref{sec:jackknife-difference-check}
specifies the correction and reports all four comparisons.

The complete design contains 104 population pairs and 312 distinct settings.
Shared settings, such as the zero-difference points in the power study,
use the same simulation results in each comparison. At 20,000 repetitions
per setting, the study samples 6,240,000 table pairs.

\section{Sampling and performance measures}

For each setting, we draw 20,000 independent multinomial count-table pairs
from its fixed \(P\) and \(Q\). In each repetition, we draw one count
table from \(P\) and another independently from \(Q\), then apply both
tests to these same two tables.

All settings use a two-sided 5\% test. The primary rejection rate is
\begin{equation}
 \frac{\#\{\text{valid p-values no larger than }0.05\}}{20{,}000}.
 \label{eq:unconditional-rate}
\end{equation}
At equal MI this estimates the false-positive rate; when MI differs it
estimates power. Invalid calculations contribute no rejection but are not
evidence for equal MI, so the valid-result frequency is reported
alongside rejection. Rates restricted to method-valid or common-valid
samples and their denominators are specified in
Appendix~\ref{app:construction}.

We report 95\% confidence intervals to show how precisely the simulations
estimate each rejection rate. These intervals use the Wilson method
\citep{wilson1927}. With 20,000 repetitions, an estimated rate of 5\%
has an interval of approximately 4.7\%--5.3\%.

When comparing Normal Wald and Expanded Welch, we calculate the confidence
interval for their rejection-rate difference using both tests' decisions
on the same simulated table pairs.
Appendix~\ref{sec:simulation-uncertainty} gives the calculations.

\medskip
\thesisminiheading{Accuracy of the MI difference and standard error.}

Bias measures whether the estimated MI difference is systematically too
high or too low. We estimate it by subtracting the true MI difference
from the mean corrected MI difference across repetitions.

The standard error should describe how much the estimated MI difference
varies between samples. Squaring each standard error and averaging across
repetitions gives the mean variance estimated by the tests. We divide
this mean estimated variance by the observed variance of the estimated MI
differences. A ratio of one indicates agreement; a larger ratio indicates
overestimated uncertainty, while a smaller ratio indicates underestimated
uncertainty.

We also check the first-order variance approximation using the known
population probabilities. We compare \(V(P)/n_P+V(Q)/n_Q\) with the
observed variance of the estimated MI differences. Using the true probabilities
separates errors in the approximation from errors caused by estimating
those probabilities. Appendix~\ref{sec:estimation-diagnostics} gives the
calculations and the rules for which repetitions are included.

\medskip
\thesisminiheading{Estimated degrees of freedom.}

We record Expanded Welch's estimated degrees of freedom in
all three studies and both supporting checks. For each setting, the median and
5th--95th percentile range summarise how these estimates vary across
valid Expanded Welch outputs. We examine how these summaries change with
population MI, sample size, table size, construction and allocation.

\medskip
\thesisminiheading{Computational cost.}

The timing comparison measures computational cost on pre-generated
\(2\times2\) through \(16\times16\) table pairs at
\(n=1{,}000\) under different skew. Sampling and population construction
are excluded from the timings. Both methods scan the cells, so their
operation counts grow linearly with the number of cells. The timing
procedure and measurement environment are given in
Appendix~\ref{sec:timing-procedure}.
